\documentclass[11pt,a4paper]{article}
\usepackage[utf8]{inputenc}
\usepackage[T1]{fontenc}
\usepackage{graphicx}
\usepackage{booktabs}
\usepackage{float}
\usepackage[section]{placeins}
\usepackage[margin=2.5cm]{geometry}
\usepackage[hidelinks,breaklinks]{hyperref}
\usepackage{amsmath}
\usepackage{enumitem}
\usepackage{setspace}
\usepackage{xurl}
\usepackage[labelfont=bf]{caption}
\usepackage{xcolor}
\newcommand{\revised}[1]{\textcolor{red}{#1}}
\onehalfspacing

\title{\LARGE\bfseries Conditional Layer Invariance in Protein Language Models:\\\revised{Position Leakage and Readout Choice Determine the Apparent Benefit of Multi-Scale Embeddings}}

\author{
\large
Yutao Guo\textsuperscript{1,2}, Zihan Zhang\textsuperscript{1,2}, Xuezhou Zhao\textsuperscript{1,2},\\[2pt]
Mengxi Chen\textsuperscript{1,2}, Dan Wu\textsuperscript{1,3*}\\[8pt]
\small
\textsuperscript{1}International College of Pharmaceutical Innovation,\\
\hspace*{1em}Soochow University, Suzhou 215123, China\\[4pt]
\textsuperscript{2}School of Pharmacy and Biomolecular Sciences, RCSI,\\
\hspace*{1em}123 St Stephen's Green, Dublin 2, D02 VN51, Ireland\\[4pt]
\textsuperscript{3}Department of Chemistry, RCSI,\\
\hspace*{1em}123 St Stephen's Green, Dublin 2, D02 YN77, Ireland\\[8pt]
*Corresponding author: \texttt{danwu@suda.edu.cn}}
\date{}

\begin{document}
\maketitle

\begin{abstract}
\textbf{Motivation:} Protein language models such as ESM-2 encode biochemical information across
33 transformer layers. Multi-layer embedding fusion is often used under the assumption
that different layers provide complementary mutation-relevant information. However,
the conditions under which this assumption holds for mutation effect prediction remain
poorly characterized.

\textbf{Results:} \revised{We benchmarked two 650M-parameter PLMs on the ProteinGym substitution
benchmark using per-layer probing under three feature arms and four train--test splits,
multi-scale pooling, logit lens scoring, layer-wise similarity analysis, and a
position-lookup control. Representational geometry changed substantially across depth
(adjacent-layer cosine $\approx 0.98$), yet for SaProt the linear CKA between layers 0
and 32 remained 0.81 despite a raw cosine of $-0.01$: the network rotates its
representations while preserving linearly decodable information. Under the standard
random split, all 33 layers again gave nearly indistinguishable predictions
($\rho \approx 0.60$), but a control regressor given only random position vectors
reproduced this level ($\rho = 0.600$) and collapsed to $0.10$--$0.14$ under
leakage-free splits (paired drop $-0.46$, Wilcoxon $p = 5.2 \times 10^{-12}$):
random-split layer invariance is largely position memorization. Under three
leakage-free splits (modulo, contiguous, position-grouped), mutation-specific
embeddings reached $\rho$ up to 0.53, the best layers were consistently late (layers
30--32), and the best-layer versus final-layer difference never exceeded
$\Delta\rho = 0.009$ (TOST equivalence at margin 0.01 in 7 of 19 tests; the remainder
inconclusive, none beyond the margin). Real Foldseek 3Di structure tokens improved
SaProt only under leakage-free evaluation ($+0.06$ to $+0.12$, all $p \leq 10^{-4}$),
with much smaller gains under the random split ($\leq +0.03$), so both the apparent
uselessness of structure tokens and the apparent layer invariance are artifacts of
position leakage. A corrected logit lens implementation (fp32, validated against the
official logits pathway) showed monotonically improving zero-shot scoring with depth
for both models, final layer best (ESM-2 $\rho = 0.443$; SaProt $\rho = 0.437$); the
previously reported non-monotonic ESM-2 curve is withdrawn. Multi-scale pooling under
the leakage-free protocol occupied a narrow band across four fusion strategies and
four readouts (20 datasets per model), with per-fold-trained attention pooling never
outperforming the last layer.}

\textbf{Availability and Implementation:} All code and processed results are available at
\url{https://github.com/schwarzeteeguo-creator/esm2-layer-invariance}. \revised{An archived
snapshot of the code and per-dataset results for this revision is deposited at Zenodo
(DOI: [AUTHOR\_INPUT\_NEEDED]).} The benchmarking pipeline, including masked-position
feature extraction, multi-scale pooling, and logit lens analysis, is implemented in Python
using PyTorch and scikit-learn.
\end{abstract}

%=============================================================================
\section{Introduction}
%=============================================================================

Protein language models (PLMs) pre-trained via masked language modeling (MLM) on
millions of protein sequences have become indispensable for mutation effect
prediction~\cite{rives2021biological,meier2021language,lin2023evolutionary}.
Among these, ESM-2~\cite{lin2023evolutionary}
is one of the most widely adopted, with a 650M-parameter architecture comprising 33
transformer layers that produce 1280-dimensional per-residue hidden states. These
embeddings have been successfully applied to diverse tasks including variant effect
prediction~\cite{frazer2021disease,hie2021learning}, protein design, and functional
annotation~\cite{elnaggar2022prottrans,brandes2022proteinbert}, establishing ESM-2 as a
foundational tool in computational biology.

The internal organization of information across transformer layers has been
extensively studied in natural language processing, where the ``BERTology''
literature~\cite{rogers2021bertology,tenney2019bert} has revealed that different layers specialize
in distinct linguistic phenomena: lower layers encode surface-level syntax, middle
layers capture semantic content, and final layers refine task-specific
predictions~\revised{\cite{tenney2019bert,rogers2021bertology}}.
This hierarchical view has motivated analogous investigations in protein language
models. A growing body of work posits that different transformer layers in PLMs
similarly encode complementary biochemical information, and that combining
representations from multiple layers should therefore improve downstream tasks.
Kumar and Jha~\cite{kumar2025} demonstrated that mid-to-late ESM-2 layers (20--33)
outperform the final layer for kinase function prediction, while early layers
(6--14) have been shown to encode local physicochemical properties such as
hydrophobicity and secondary structure propensity~\cite{detlefsen2022}.
Alsamkary, Elshaffei et al.~\cite{elshaffei2025} proposed hierarchical pooling and
pooled attention architectures for fusing representations across multiple PLMs and
reported up to a 12\% Spearman improvement for protein--protein interaction
prediction. More broadly, the strategy of extracting and fusing embeddings from
multiple transformer layers has become common practice in PLM-based
pipelines~\cite{rao2019tape}, often under the implicit assumption that each layer
contributes non-redundant information.

These claims, however, rest on an assumption that has not been systematically tested:
that different ESM-2 layers genuinely encode \emph{complementary} information for
mutation effect prediction, rather than merely offering different views of the same
underlying signal. If layers are functionally similar under a given readout and
evaluation protocol, then multi-scale fusion may add cost without improving prediction.
If layer complementarity appears only under specific readouts or split protocols,
then layer selection should be treated as a conditional design choice rather than
a universal rule. Resolving this question carries practical
consequences: the computational cost of extracting and storing embeddings from
multiple layers is substantial for large-scale screening efforts, and the
proliferation of ad-hoc layer-selection heuristics complicates reproducibility.

To resolve this question, we conducted a large-scale benchmark on ProteinGym~\cite{notin2024proteingym},
a comprehensive collection of deep mutational scanning (DMS) experiments spanning
diverse protein families, assay types, and organismal contexts. We employed
three independent experimental paradigms (Figure~\ref{fig:concept}): (1) per-layer supervised probing via Ridge
regression, (2) multi-scale embedding pooling with controlled feature sets, and
(3) logit lens zero-shot scoring. In parallel, we quantified layer-wise representational
similarity through cosine similarity analysis. Together, these approaches test when
PLM layers behave as practically interchangeable representations for mutation effect
prediction and when readout choice or evaluation protocol exposes layer-specific
differences. \revised{We report both the artifacts and the leakage-free re-analysis:
the random-split layer profile is largely position memorization, while leakage-free
evaluation reveals small, consistent late-layer advantages and a measurable benefit of
structure tokens that the standard protocol cannot see.}

\begin{figure}[ht]
\centering
\includegraphics[width=\textwidth]{figures/output/fig0_concept.pdf}
\caption{\textbf{Experimental design overview.}
The three complementary experimental paradigms employed in this study:
(\textbf{a}) Per-layer probing: Ridge regression on hidden states extracted from each
of the 33 transformer layers \revised{under three feature arms (masked wild-type,
unmasked wild-type, and mutant) and four train--test splits (random, modulo,
contiguous, and position-grouped), quantifying the predictive information content per
layer together with a random-position lookup control}.
(\textbf{b}) Multi-scale pooling: four embedding fusion strategies (Last Layer,
Mean L20--33, Concat 5L, Attention-weighted\revised{, with attention weights trained
within each cross-validation fold}) tested across four readout models
(Ridge, Random Forest, MLP, LightGBM) to assess whether combining layers improves
prediction.
(\textbf{c}) Logit lens: applying the final-layer LM head to intermediate hidden
states for zero-shot mutation scoring, directly measuring how token predictions
are refined across depth.
(\textbf{d}) Layer-wise cosine similarity analysis: quantifying representational
divergence independent of predictive performance.
Together, these paradigms provide converging evidence on whether different
transformer layers encode complementary or redundant information for mutation
effect prediction.}
\label{fig:concept}
\end{figure}

%=============================================================================
\section{Methods}
%=============================================================================

\subsection{Models and Data}

We used SaProt\_650M\_AF2~\cite{su2024saprot}, a structure-aware variant of ESM-2 650M
that extends the vocabulary to 446 structure-informed tokens (21 amino acids $\times$ 21
Foldseek 3Di states~\cite{vankempen2024foldseek}), as our primary model. For validation,
we also tested standard ESM-2 650M~\cite{lin2023evolutionary} in both the probing and
logit lens experiments, enabling direct comparison between structure-aware and
sequence-only representations. All experiments used ProteinGym's substitution DMS
benchmark~\cite{notin2024proteingym}, comprising 63 datasets with experimentally measured
fitness effects for single and multiple amino acid substitutions. \revised{All 63 datasets
were used for the probing and logit lens experiments ($n = 63$ for the random, modulo, and
position-grouped splits; six datasets whose mutated positions form a single contiguous
block have undefined Spearman correlations under the contiguous split, which is therefore
reported on the remaining 57).} The pooling benchmark was run on a computationally tractable subset of
20 datasets because it crossed four pooling strategies with four readout models and
multiple cross-validation folds. We therefore treat the pooling analysis as a
controlled validation of the probing result rather than as a full ProteinGym-wide
estimate. \revised{The subset comprises the first 20 datasets in alphabetical order,
fixed identically for both models.} Per-dataset mutation counts range from under 100 to over 200,000.
\revised{For SaProt, we additionally prepared real Foldseek 3Di structure tokens by
assigning each wild-type residue the 3Di state of its aligned residue in the
AlphaFold Protein Structure Database entry for that sequence (structures available for
62 of 63 datasets; two aligned structures were truncated at the C terminus and the
missing residues were padded with the Foldseek gap state; the influenza nucleoprotein
dataset has no AFDB entry and was evaluated with structure-masked tokens). SaProt
results are reported for both structure-masked and real-3Di inputs.}

\subsection{Experiment 1: Per-Layer Probing}

\revised{For each DMS dataset we extracted per-layer hidden states under three feature
arms. The \textbf{masked\_wt} arm replaces each mutated position in the wild-type
sequence with the mask token (``M\#'' for SaProt, \texttt{<mask>} for ESM-2) and reads
the hidden state at that position, exactly as in the submitted version. The
\textbf{unmasked\_wt} arm reads the wild-type token at the mutated position without
masking. The \textbf{mutant} arm feeds the full mutant sequence and reads the hidden
state at the mutated position; it is the only arm whose representation differs between
two mutations at the same position, and therefore the only one that can carry
mutation-specific information beyond position identity. All mutations of a variant
were masked jointly in the masked\_wt arm. Both SaProt\_650M\_AF2 (structure-masked
and real-3Di inputs) and standard ESM-2 650M were evaluated under identical conditions.}

\revised{Each arm was evaluated under four train--test splits, all with five
cross-validation folds (four for the single four-position dataset under the
position-grouped split) and with all rows of a multi-mutant variant assigned to the
same fold: the standard \textbf{random} split over mutations; the \textbf{modulo}
split, which assigns each residue position to fold $(\text{position} - 1) \bmod 5$
following the official ProteinGym protocol; the \textbf{contiguous} split, which
divides the sequence into five contiguous segments and holds out one segment per
fold; and a \textbf{position-grouped} split (GroupKFold over residue positions, so
that all mutations at a given position share a fold). The latter three are
leakage-free with respect to positions. For each layer $l \in \{0,\dots,32\}$ and
feature configuration (combined: the hidden state concatenated with 42 auxiliary
features, namely one-hot encodings of wild-type and mutant amino acids, the BLOSUM62
substitution score, and a normalized Grantham distance proxy; embedding\_only; and
auxiliary\_only), we trained Ridge regression with the regularization strength
selected per dataset by generalized cross-validation~\cite{golub1979generalized}
(grid $10^{-2}$ to $10^{3}$), standardization fit inside each training fold, and
reported Spearman rank correlation $\rho$. Datasets were subsampled to at most 10,000
rows before feature construction and 5,000 rows inside the pipeline (11 of 63 datasets
were capped), with a fixed seed shared across arms, splits, and layers. The single
dataset with fewer than five mutated positions (four positions) used four
position-grouped folds.}

\revised{To quantify how much probing performance is attributable to position identity
alone, we added a \textbf{random-position control}: each mutation is represented by a
fixed random vector indexed by its residue position (33 independent draws), probed with
the identical Ridge protocol with and without auxiliary features. This control contains
no sequence information; any performance it recovers measures position memorization.}

\subsection{Experiment 2: Multi-Scale Pooling}

\revised{We compared four embedding strategies using fixed auxiliary features,
evaluated on the \textbf{mutant} arm under the modulo split as the primary
leakage-free protocol: (1) \textbf{Last Layer}: the standard approach using only the
final transformer layer;
(2) \textbf{Mean L20--33}: mean pooling across layers 20--33, motivated by
Kumar and Jha~\cite{kumar2025};
(3) \textbf{Concat 5L}: concatenation of layers 6, 14, 20, 26, and 32 (0-indexed,
including the final layer), spanning the full representational depth; and
(4) \textbf{Attention-Weighted}: softmax-weighted pooling across all 33 layers, with
the weights trained within each cross-validation fold.}

\revised{For each strategy, we evaluated Ridge regression (regularized by generalized
cross-validation as above), Random Forest (50 estimators,
max depth 10), MLP (two hidden layers of 256 units, ReLU), and LightGBM (100 estimators,
max depth 7) with 3-fold cross-validation for datasets exceeding 3,000 mutations
(5-fold otherwise). Standardization was fit inside training folds, and the auxiliary
features were identical across strategies, isolating the effect of embedding choice.}

\subsection{Experiment 3: Logit Lens}

\revised{We applied the model's final-layer LM head to the hidden states of each
intermediate layer, producing pseudo-logits for zero-shot mutation scoring. All
computations ran in fp32, and for every dataset the final-layer pseudo-logits were
validated against the model's official logits pathway (maximum $|\Delta \log P| <
10^{-3}$; all 63 datasets passed for both models). For standard ESM-2, we used the
conventional log-ratio, $\log P(\text{mutant}) - \log P(\text{wild-type})$. For SaProt
with structure-masked input tokens ($\text{AA}\#$, where $\#$ denotes masked
structure), we marginalized over all 21 structure-token variants per amino acid; with
real 3Di inputs, the wild-type structure token at each position was retained. We
reported per-layer Spearman $\rho$ between these zero-shot scores and experimental
DMS fitness values across all 63 datasets. The lens curves published in the submitted
version were produced by an fp16 pipeline whose numerical error invalidated the ESM-2
curve; they are withdrawn and replaced by these validated results
(cf.\ the tuned lens literature~\cite{belrose2023eliciting}).}

\subsection{Layer Similarity Analysis}

\revised{We computed the $33 \times 33$ cosine similarity matrix of per-layer hidden states,
averaged over 50 randomly selected residue positions across 10 datasets. Because raw
cosine is not invariant to rotations of the coordinate system, we additionally computed
linear CKA~\cite{kornblith2019similarity} on the same hidden-state matrices. The unit
of analysis for similarity-versus-distance statistics is the set of unique layer pairs
$(i, j)$ with $i < j$, diagonal excluded. A high CKA combined with a near-zero cosine
indicates a change of basis that preserves linearly decodable information.}

\subsection{Statistical Evaluation}

We used three complementary metrics: (1) \textbf{invariance score}, defined as
the standard deviation of per-layer mean Spearman $\rho$ (smaller indicates greater invariance);
\revised{(2) \textbf{equivalence testing}: for each model, arm, and leakage-free split,
we tested whether the mean per-dataset difference between the best layer and the final
layer falls within a margin of $\Delta\rho = 0.01$ using two one-sided tests (TOST) on
paired differences, with the best layer fixed as the argmax of the layer-mean curve
(19 tests in total), each reported with its 90\% bootstrap confidence interval and
Cohen's $d_z$; and (3) \textbf{paired Wilcoxon signed-rank tests} for split contrasts
(random versus leakage-free) and for real-3Di versus structure-masked SaProt inputs.}
We use the term \emph{practical layer invariance} operationally rather than as a
formal proof of mathematical equivalence. A layer set was considered practically
invariant when per-layer variation was small relative to between-dataset variation,
when the best-layer versus final-layer difference had negligible effect size, and
when no layer comparison remained significant after false-discovery-rate correction.
This framing is intentionally conservative: non-significant hypothesis tests alone
are not interpreted as proof that two layers are identical.

%=============================================================================
\section{Results}
%=============================================================================

\subsection{Representations Change Dramatically While Predictive Performance Is Preserved}

We first asked whether different transformer layers encode representations that are
merely different views of the same underlying information, or whether they capture
genuinely complementary features. To answer this, we computed the $33 \times 33$ layer-wise
cosine similarity matrix of per-layer hidden states, averaged over 50 randomly selected
residue positions across 10 datasets (Figure~\ref{fig:similarity}).

\begin{figure}[ht]
\centering
\includegraphics[width=0.7\textwidth]{figures/output/fig2_similarity_heatmap.png}
\caption{\textbf{Layer-wise cosine similarity matrix (33$\times$33).}
Averaged over 10 datasets and 50 random residue positions per dataset.
Adjacent layers are highly similar ($\cos=0.98$); \revised{mean similarity decays
monotonically with distance and collapses to $-0.01$ only between the most distant
layers (Layer~0 vs.~32). Linear CKA between layers 0 and 32 remains 0.81 for SaProt,
indicating a rotation that preserves linearly decodable information rather than
information loss.} Representations change continuously even when supervised predictive
performance is nearly unchanged under random-split probing.}
\label{fig:similarity}
\end{figure}

The matrix revealed a smooth representational trajectory across layers. Adjacent layers
were highly similar ($\cos = 0.979 \pm 0.011$), and similarity decayed approximately
linearly with layer distance (Pearson $r = -0.50$). \revised{The decay was smooth and
monotonic, with no intermediate trough: for SaProt the mean cosine by distance fell
from 0.97 ($d=2$) through 0.80 ($d=20$) to 0.50 ($d=30$), collapsing to $-0.009$
between the most distant layers (Layer~0 vs.\ Layer~32). ESM-2 decayed more gently and
remained positive throughout (0.93 at $d=2$, 0.56 at $d=32$). Crucially, near-zero raw
cosine does not imply lost information: the linear CKA between SaProt layers 0 and 32
was 0.808 while their cosine was $-0.009$ (ESM-2: CKA 0.554, cosine 0.558; adjacent-layer
CKA 0.97 for both models). The depth-wise transformation is therefore well described as
a rotation of the representation that preserves linearly decodable information, which
explains how supervised performance can remain flat across depth while raw cosine
collapses (Figure~\ref{fig:fig4}d).} This large
rotation of the representational geometry raises a sharp question: \emph{does this
representational transformation reflect complementary information encoded at different
depths, or is mutation-relevant predictive performance preserved under this transformation?}

To answer this, we conducted three converging experimental paradigms, each designed to
interrogate the information content of individual layers from a different angle.

\subsection{\revised{Random-Split Layer Invariance Is Largely Position Memorization}}

We first asked \revised{whether the near-flat random-split layer profile of the submitted
version reflects genuine representational redundancy. Under the unified pipeline it
replicates exactly: for the masked\_wt arm, per-layer mean $\rho$ spans
$0.603$--$0.606$ for both models with invariance scores $\sigma \approx 5 \times
10^{-4}$ to $1.2 \times 10^{-3}$ (Table~\ref{tab:probing}). A Ridge regressor given
\emph{random vectors indexed only by mutation position} plus the same auxiliary
features recovered $\rho = 0.600$ under the random split (Table~\ref{tab:control}),
matching the full masked-position PLM probing level. Under leakage-free splits the
same control collapsed to $\rho = 0.138$ (modulo), $0.096$ (contiguous, $n=57$), and
$0.135$ (position-grouped), a paired drop of $-0.46$ (random versus modulo, Wilcoxon
$p = 5.2 \times 10^{-12}$); the random vectors alone carried $0.557$ under the random
split and approximately zero elsewhere, and the auxiliary features alone fell from
$0.339$ to $0.19$--$0.23$. Because every mutation at the same position shares one
masked-position embedding, the random-split task is solvable by a position lookup
table; the flat layer profile is therefore largely an artifact of position
memorization rather than evidence that mutation information is distributed uniformly
across depth.}

\begin{table}[ht]
\centering
\caption{\revised{Random-position lookup control (63 datasets; contiguous $n = 57$).
Mean Spearman $\rho$ of Ridge regressors given no sequence information. The control
reproduces the random-split PLM probing level and collapses under leakage-free
splits.}}
\label{tab:control}
\begin{tabular}{lcccc}
\toprule
Features & Random & Modulo & Contiguous & Position-grouped \\
\midrule
Random position vector + aux & 0.600 & 0.138 & 0.096 & 0.135 \\
Random position vector only  & 0.557 & $-0.008$ & $-0.056$ & $-0.010$ \\
Auxiliary only               & 0.339 & 0.230 & 0.188 & 0.228 \\
\bottomrule
\end{tabular}
\end{table}

\revised{Table~\ref{tab:probing} and Figure~\ref{fig:probing}a report the unified
re-probing of all three arms,
both models, and all four splits ($n = 63$; contiguous $n = 57$). Two patterns are
stable across every leakage-free column. First, the best layer sits in a narrow late
band: layers 30--32 are best in 26 of 27 model--arm--split combinations (layer 27 in
the remaining one). Second, the best layer is at most marginally better than the
final layer: the mean per-dataset difference never exceeded $\Delta\rho = 0.009$,
TOST at the 0.01 margin declared equivalence in 7 of 19 leakage-free tests, and the
remaining 12 were inconclusive (90\% confidence intervals overlapping the margin); no
test indicated a difference beyond the margin (Supplementary Table S4). The
invariance score rose from $\sigma \approx 10^{-3}$ (masked\_wt, random) to
$4$--$7 \times 10^{-2}$ under leakage-free evaluation, yet remained well below the
cross-dataset standard deviation ($\approx 0.17$): performance varies far more
\emph{across} datasets than \emph{across layers}, but under leakage-free evaluation
the residual layer variation has a consistent direction, favoring the final layers.
The mutant arm exceeded the masked arms in every split (e.g., 0.681 versus 0.606
under random, 0.483 versus 0.459 under modulo for ESM-2), confirming that
mutation-specific signal exists and is recoverable once the representation actually
depends on the mutation.}

\begin{table}[ht]
\centering
\caption{\revised{Unified per-layer probing, combined features (mean Spearman
$\rho$; best layer in parentheses, followed by the final-layer value when the best
layer is not layer 32). $n = 63$; contiguous $n = 57$. Full 33-layer curves in
Supplementary Table S2; per-dataset results in Supplementary Table S3.}}
\label{tab:probing}
\begin{tabular}{llcccc}
\toprule
Model & Arm & Random & Modulo & Contiguous & Pos.-grouped \\
\midrule
ESM-2 & masked\_wt   & 0.606 (32) & 0.459 (32) & 0.370 (32) & 0.450 (32) \\
ESM-2 & unmasked\_wt & 0.606 (32) & 0.470 (32) & 0.388 (32) & 0.461 (32) \\
ESM-2 & mutant       & 0.681 (32) & 0.483 (31) $\rightarrow$ 0.479 & 0.365 (31) $\rightarrow$ 0.360 & 0.477 (31) $\rightarrow$ 0.475 \\
\midrule
SaProt (mask) & masked\_wt   & 0.603 (31) $\rightarrow$ 0.603 & 0.391 (30) $\rightarrow$ 0.386 & 0.282 (31) $\rightarrow$ 0.274 & 0.385 (31) $\rightarrow$ 0.382 \\
SaProt (mask) & unmasked\_wt & 0.604 (32) & 0.412 (30) $\rightarrow$ 0.406 & 0.324 (31) $\rightarrow$ 0.320 & 0.403 (27) $\rightarrow$ 0.400 \\
SaProt (mask) & mutant       & 0.652 (23) $\rightarrow$ 0.651 & 0.465 (30) $\rightarrow$ 0.460 & 0.323 (32) & 0.461 (30) $\rightarrow$ 0.458 \\
\midrule
SaProt (3Di) & masked\_wt   & 0.606 (32) & 0.461 (31) $\rightarrow$ 0.460 & 0.394 (31) $\rightarrow$ 0.394 & 0.454 (31) $\rightarrow$ 0.453 \\
SaProt (3Di) & unmasked\_wt & 0.606 (27) $\rightarrow$ 0.605 & 0.473 (31) $\rightarrow$ 0.466 & 0.415 (31) $\rightarrow$ 0.406 & 0.467 (31) $\rightarrow$ 0.459 \\
SaProt (3Di) & mutant       & 0.684 (26) $\rightarrow$ 0.683 & 0.531 (31) $\rightarrow$ 0.525 & 0.451 (31) $\rightarrow$ 0.445 & 0.532 (31) $\rightarrow$ 0.524 \\
\bottomrule
\end{tabular}
\end{table}

\subsection{\revised{Multi-Scale Pooling Shows No Practical Gain Under a Leakage-Free Protocol}}

\revised{We re-ran the pooling benchmark with the unified pipeline (mutant arm, modulo
split, GCV-regularized Ridge, fold-internal standardization, attention weights trained
per fold) on the 20-dataset design subset for each model
(Table~\ref{tab:pooling}; Figure~\ref{fig:probing}b).
For SaProt (structure-masked inputs), all four strategies occupied a narrow band for
every readout (Ridge spread across strategies 0.005; LightGBM 0.014). For ESM-2,
spreads were likewise small (Ridge 0.041; LightGBM 0.050), and the best single
configuration was the last layer with LightGBM ($\rho = 0.508$), exceeding every
multi-scale strategy under every readout. Per-fold-trained attention pooling was
never the best strategy for either model.}

\begin{table}[ht]
\centering
\caption{\revised{Multi-scale pooling under the unified leakage-free pipeline (mutant
arm, modulo split, combined features; 20 datasets per model; mean $\rho \pm$ sd).
Spreads across strategies: ESM-2 Ridge 0.041, RF 0.057, MLP 0.045, LightGBM 0.050;
SaProt Ridge 0.005, RF 0.010, MLP 0.039, LightGBM 0.014.}}
\label{tab:pooling}
\begin{tabular}{lcccc}
\toprule
\textbf{ESM-2 650M} & Ridge & RF & MLP & LightGBM \\
\midrule
Last Layer (baseline) & $0.463 \pm 0.174$ & $0.476 \pm 0.173$ & $0.480 \pm 0.198$ & $\mathbf{0.508} \pm 0.189$ \\
Mean L20--33          & $0.473 \pm 0.180$ & $0.420 \pm 0.165$ & $0.436 \pm 0.179$ & $0.457 \pm 0.174$ \\
Concat 5L             & $0.481 \pm 0.183$ & $0.453 \pm 0.181$ & $0.435 \pm 0.170$ & $0.495 \pm 0.190$ \\
Attention-weighted    & $0.439 \pm 0.201$ & $0.419 \pm 0.209$ & $0.440 \pm 0.209$ & $0.463 \pm 0.208$ \\
\midrule
\textbf{SaProt\_650M\_AF2} & Ridge & RF & MLP & LightGBM \\
\midrule
Last Layer (baseline) & $0.432 \pm 0.163$ & $0.415 \pm 0.174$ & $0.445 \pm 0.168$ & $\mathbf{0.448} \pm 0.178$ \\
Mean L20--33          & $0.433 \pm 0.166$ & $0.415 \pm 0.166$ & $0.429 \pm 0.166$ & $0.439 \pm 0.182$ \\
Concat 5L             & $0.434 \pm 0.185$ & $0.407 \pm 0.173$ & $0.406 \pm 0.176$ & $0.436 \pm 0.183$ \\
Attention-weighted    & $0.429 \pm 0.164$ & $0.405 \pm 0.166$ & $0.426 \pm 0.164$ & $0.433 \pm 0.177$ \\
\bottomrule
\end{tabular}
\end{table}

\subsubsection{\revised{The Previously Reported ESM-2 Asymmetry Was a Cross-Protocol Artifact}}

\revised{The large ESM-2 asymmetry in the submitted version (last-layer Ridge
$\rho = 0.386$ versus Concat~5L $\rho = 0.639$) does not survive protocol
unification. It arose from comparing feature constructions that differed in masking
state, regularization, and split protocol. Under the unified leakage-free pipeline
the corresponding contrast is $0.463$ versus $0.481$ ($+0.018$), inside the narrow
band described above. The MLP readout, previously reported at $\rho \approx 0.33$
across strategies, reflected an under-trained surrogate configuration in the
submitted pipeline (one hidden layer of 128 units, 50 iterations) rather than the
described architecture; with the described two-layer architecture it performs on par
with Ridge here. We therefore withdraw the claim that standard ESM-2 exhibits a
large readout-specific layer dependence; residual readout interactions are an order
of magnitude smaller under the unified protocol, and no strategy consistently beats
the last layer.}

\begin{figure}[ht]
\centering
\includegraphics[width=\textwidth]{figures/output/fig4_summary.png}
\caption{\textbf{\revised{Layer invariance and its boundary conditions across three experimental paradigms.}}
(\textbf{a}) \revised{Multi-scale pooling under the unified leakage-free pipeline (mutant arm,
modulo split, 20 datasets per model): both models occupy a narrow band across strategies
(ESM-2 Ridge 0.439--0.481; SaProt Ridge 0.429--0.434), and last-layer LightGBM is the best
single configuration for both models.}
(\textbf{b}) \revised{Logit lens with ESM-2 ($n = 63$): zero-shot $\rho$ rises with depth to
0.443 at the final layer.} (\textbf{c}) \revised{Logit lens with SaProt real-3Di inputs
($n = 63$): zero-shot $\rho$ rises from negative early-layer values to 0.437 at the final layer.}
(\textbf{d}) \revised{Cosine similarity versus layer distance (unique layer pairs, density-colored
scatter): SaProt decays monotonically and reaches $-0.01$ only at the largest distances, while
ESM-2 remains positive; linear CKA between the most distant SaProt layers remains 0.81,
indicating rotation with preserved linear information rather than loss.}}
\label{fig:fig4}
\end{figure}

\begin{figure}[ht]
\centering
\includegraphics[width=\textwidth]{figures/output/fig3_probing_pooling.png}
\caption{\textbf{\revised{Per-layer probing and multi-scale pooling under the unified protocol.}}
(\textbf{a}) \revised{33-layer probing (combined features): under the random split the
masked-arm profile is flat for both models ($\sigma \approx 10^{-3}$); under leakage-free
splits performance is lower and the best layers are consistently late (layers 30--32,
$\Delta\rho \leq 0.009$ versus the final layer).}
(\textbf{b}) \revised{Pooling (mutant arm, modulo split, 20 datasets per model): ESM-2 spread
across strategies 0.041 (Ridge) to 0.057 (RF); SaProt spread 0.005 (Ridge) to 0.039 (MLP);
per-fold-trained attention pooling is never best.}}
\label{fig:probing}
\end{figure}

\subsection{Logit Lens \revised{Confirms Late-Layer Information Accumulation}}

\revised{The re-implemented lens (fp32, per-dataset validation against the official
logits pathway, $n = 63$) produced monotonically improving zero-shot scores with
depth for both models (Figure~\ref{fig:logit_lens}; Figure~\ref{fig:fig4}b,c). ESM-2 rose from $\rho = 0.031$
(layer 0) to $\rho = 0.443$ (layer 32), with only shallow fluctuations in early
layers; SaProt with real 3Di tokens rose from $-0.077$ (layer 0) through
approximately zero around layer 20 to $\rho = 0.437$ (layer 32). Both curves attain
their maximum at the final layer, at magnitudes consistent with published zero-shot
PLM scoring. The submitted version instead reported a non-monotonic ESM-2 curve
peaking at layer 1 ($\rho = 0.056$) and ending at $-0.021$; that curve was produced
by an fp16 pipeline and is withdrawn. The corrected result removes the apparent
model contrast and agrees with the leakage-free probing: mutation-relevant
information accumulates toward the final layers. Zero-shot absolute levels remain
below supervised probing (0.44 versus up to 0.53), so the lens is best read as a
consistency check on layer-wise trends rather than a standalone benchmark.}

\begin{figure}[ht]
\centering
\includegraphics[width=0.85\textwidth]{figures/output/fig1_logit_lens.png}
\caption{\textbf{Logit lens: zero-shot mutation scoring by layer ($n = 63$).}
\revised{Both models improve with depth and attain their maximum at the final layer:
ESM-2 (orange) rises from $\rho = 0.03$ (layer 0) to $\rho = 0.443$ (layer 32);
SaProt with real 3Di tokens (blue) rises from negative early-layer values to
$\rho = 0.437$ (layer 32). Shaded bands indicate $\pm 1$ standard deviation across
datasets. All computations in fp32 and validated against the official logits
pathway; the previously reported non-monotonic ESM-2 curve is withdrawn.}}
\label{fig:logit_lens}
\end{figure}

\subsection{Boundary Conditions}

The results above provide evidence for practical layer invariance under common
random-split supervised benchmarking conditions. \revised{We now examine the two
factors that moderate this conclusion: evaluation protocol and input
representation.}

\subsubsection{\revised{Leakage-Free Splits: What Is Lost and What Remains}}

\revised{The control and main-table results above already localize the leakage: the
random-position lookup table reproduces the random-split level ($0.600$ versus
$0.603$--$0.606$ for full probing) and loses it entirely under leakage-free splits
($0.10$--$0.14$), while the random vectors alone carry $0.557$ under the random
split and approximately zero elsewhere (Tables~\ref{tab:control}
and~\ref{tab:probing}; Figure~\ref{fig:position_level}). Among the leakage-free
splits, modulo and position-grouped give similar levels (e.g., SaProt real-3Di
mutant arm: $0.531$ and $0.532$) while contiguous is uniformly hardest ($0.451$),
consistent with sequence-contiguous segments correlating with local structural
context and mutational scope. Random mutation splitting, the de facto standard in DMS
benchmarking, therefore inflates absolute performance estimates by roughly
$0.15$--$0.20$ $\rho$ in these arms and masks both layer differences and
structure-token contributions. We recommend that future DMS benchmarks report at
least one leakage-free split alongside the standard random split, analogous to the
time-series community's practice of reporting both random and temporal
cross-validation~\cite{roberts2017crossvalidation}.}

\begin{figure}[ht]
\centering
\includegraphics[width=\textwidth]{figures/output/fig5_position_level.png}
\caption{\textbf{\revised{Random versus leakage-free split probing (SaProt, 63 datasets).}}
(\textbf{a}) \revised{Random mutation split: flat per-layer profile
($\sigma \approx 10^{-3}$), reproduced by the random-position lookup control
($\rho = 0.600$).} (\textbf{b})
\revised{Leakage-free splits (modulo, position-grouped, contiguous): performance
drops to $0.27$--$0.53$ depending on arm and model, the invariance score rises to
$4$--$7 \times 10^{-2}$, and the best layers shift consistently to layers 30--32.}}
\label{fig:position_level}
\end{figure}

\revised{Together, these boundary conditions refine the central finding: apparent
layer invariance under random splitting is dominated by position memorization,
leakage-free evaluation reveals small but consistent late-layer advantages, and the
benefit of structure tokens is itself conditional on the evaluation protocol.}

\subsection{\revised{Structure Tokens Help Only Under Leakage-Free Evaluation}}

\revised{Because SaProt can be run with real 3Di tokens or with structure-masked
tokens, the unified protocol allows a paired test, dataset by dataset and layer by
layer, of what structure tokens actually contribute (Table~\ref{tab:3di}). Under the
random split the two input modes were nearly identical (paired differences at layer
32: $+0.003$ masked\_wt, $+0.002$ unmasked\_wt, $+0.032$ mutant). Under every
leakage-free split the real-3Di inputs were consistently superior, with mean paired
advantages of $+0.059$ to $+0.122$ across arms and splits, all with Wilcoxon
$p \leq 10^{-4}$. Structure tokens therefore carry genuine mutation-relevant
information, but that contribution becomes visible only when position memorization
is removed from the evaluation. This conditional benefit also retroactively
explains why structure-aware inputs appeared useless in random-split comparisons:
the position channel, available to both input modes, saturates the task.}

\begin{table}[ht]
\centering
\caption{\revised{Real-3Di versus structure-masked SaProt inputs (paired per-dataset
differences at layer 32, combined features; Wilcoxon signed-rank $p$). The benefit of
structure tokens is small and mostly non-significant under the random split and
large and highly significant under all leakage-free splits.}}
\label{tab:3di}
\begin{tabular}{lcc|ccc}
\toprule
& \multicolumn{2}{c|}{\textbf{Random}} & \multicolumn{3}{c}{\textbf{Leakage-free}} \\
Arm & $\Delta\rho$ & $p$ & Modulo & Contiguous & Pos.-grouped \\
\midrule
masked\_wt   & $+0.003$ & $2.2 \times 10^{-2}$ & $+0.073$ ($3.5 \times 10^{-7}$) & $+0.120$ ($8.7 \times 10^{-8}$) & $+0.071$ ($8.4 \times 10^{-7}$) \\
unmasked\_wt & $+0.002$ & $9.1 \times 10^{-1}$ & $+0.060$ ($4.6 \times 10^{-6}$) & $+0.086$ ($2.1 \times 10^{-6}$) & $+0.059$ ($1.2 \times 10^{-5}$) \\
mutant       & $+0.032$ & $2.5 \times 10^{-6}$ & $+0.064$ ($6.9 \times 10^{-5}$) & $+0.122$ ($6.5 \times 10^{-7}$) & $+0.067$ ($4.6 \times 10^{-5}$) \\
\bottomrule
\end{tabular}
\end{table}

%=============================================================================
\section{Discussion}
%=============================================================================

\revised{The central result of this study is conditional rather than absolute. Under
the standard random-split protocol, the 33 transformer layers of SaProt and ESM-2 show
near-flat mutation effect prediction performance, but the position-lookup control shows
that this flatness is largely position memorization rather than distributed mutation
information. Under leakage-free evaluation, genuine mutation-specific performance is
lower ($\rho$ up to 0.53), small but consistent late-layer advantages emerge (best
layers 30--32, $\Delta\rho \leq 0.009$), and structure tokens deliver a measurable
benefit ($+0.06$ to $+0.12$) that is invisible under random splitting. Multi-scale
pooling provides no practical gain for either model under the unified pipeline. The
practical conclusion is therefore not that multi-scale embeddings are never useful,
but that their apparent benefit is determined by the evaluation protocol and the
readout, and that leakage-free evaluation is required to see it.}

\subsection{Why Might Layers Appear Invariant?}

We propose that this practical invariance may be a consequence of the MLM pre-training
objective. During pre-training, all layers receive gradients through the same
cross-entropy loss computed at the final LM head. While this setup allows layers
to develop distinct representational formats (as the cosine similarity analysis
shows), it may not strongly incentivize each depth to specialize in fundamentally
different mutation-relevant factors. Consequently, a supervised readout trained on
individual layers can recover very similar performance under random-split evaluation
\revised{, where position memorization dominates the task}.

\revised{The cosine similarity between layers decays monotonically with distance and,
for SaProt, collapses to approximately zero only between the most distant layers,
while the linear CKA between the same layers remains 0.81. This combination, a
near-rotation that preserves linearly decodable information, is consistent with the
effect of residual connections~\cite{he2016deep}: because residual streams accumulate
representations from all preceding layers, late layers retain a linearly accessible
trace of early-layer content even as their raw coordinates rotate. The U-shaped
recovery described in the submitted version is not present in the recomputed curves
and is withdrawn.}

This interpretation is consistent with the logit lens
hypothesis~\cite{nostalgebraist2020,belrose2023eliciting},
which posits that intermediate layers in transformer language models progressively
refine token predictions toward the final output distribution, rather than
constructing qualitatively different representations. \revised{The corrected lens
curves support this directly for both models: zero-shot scoring improves with depth
and peaks at the final layer.} The near-orthogonality of the most distant SaProt
layers ($\cos \approx -0.01$) \revised{alongside their high linear CKA (0.81)}
combined with near-flat supervised probing
performance suggests that mutation-relevant signal may remain accessible while
representations move through successive coordinate frames.

Notably, this pattern contrasts with findings from NLP transformers, where
probing classifiers consistently recover different linguistic features from
different layer depths~\cite{tenney2019bert,rogers2021bertology}. In BERT, for example,
surface syntactic features peak in early layers while semantic features dominate
in middle-to-late layers. The relative flatness of the PLM probing curve may
reflect the fundamentally different nature of the learning signal in protein
versus natural language: whereas language contains hierarchical structure at
multiple timescales (morphology, syntax, semantics, discourse), the MLM task on
protein sequences operates on a single level of abstraction (amino acid identity),
potentially eliminating the pressure for layer-wise specialization.

\subsubsection{\revised{What Separates SaProt from Standard ESM-2?}}

\revised{Under the unified protocol the two models differ mainly in two measured
ways: how much their final representations retain early-layer geometry (ESM-2
distant-layer cosine 0.56 with CKA 0.55; SaProt cosine $-0.01$ with CKA 0.81), and
the benefit conferred by structure tokens. The submitted version interpreted the
pooling contrast between the models as structure-aware training distributing
information more uniformly across layers; that contrast dissolved under protocol
unification, and we no longer draw that inference. The surviving, directly measured
difference is the conditional structure-token benefit: real 3Di inputs improve
leakage-free probing by $+0.06$ to $+0.12$ while leaving random-split performance
essentially unchanged. Structure information is genuinely present and genuinely
useful, but only leakage-free evaluation can see it.}

\revised{This framing carries a testable prediction: fine-tuning standard ESM-2
with structural objectives (e.g., contact prediction or fold classification)
should reproduce a similar leakage-free advantage over the sequence-only baseline.
This remains a mechanistic hypothesis rather than a direct result of the present
experiments.}

\subsection{Reconciling with Prior Layer-Selection Studies}

Our findings do not contradict the observation that different layers may be
optimal for different \emph{task types}. Kumar and Jha~\cite{kumar2025}
demonstrated that mid-to-late ESM-2 layers (20--33) outperform the final layer
for kinase function prediction, achieving 32\% higher clustering quality and
superior classification accuracy. This is compatible with our narrower conclusion:
different downstream tasks may induce different optimal readout directions, while
mutation effect prediction under random-split probing shows little layer dependence.
The key distinction is between \emph{task-level} specialization and
\emph{within-task} layer complementarity under a specified readout and evaluation
protocol.

Similarly, the 12\% Spearman improvement reported by Alsamkary, Elshaffei
et al.~\cite{elshaffei2025} for attention-based fusion architectures was
observed in the context of multi-chain protein--protein interaction prediction,
a task that differs substantially from per-residue mutation effect scoring.
The benefit of sophisticated pooling may therefore be task-dependent rather than
inherent to multi-scale representation fusion. For mutation effect prediction,
our results suggest that such improvements should be interpreted in light of
model type, readout choice, and split protocol.

\subsection{Practical Implications}

For researchers using protein language model embeddings for mutation effect
prediction, our results support two practical recommendations that depend on
the model variant, readout model, and evaluation protocol:

\begin{enumerate}
    \item \textbf{Structure-aware models (SaProt, ESM-2 fine-tuned with structural
    objectives):} The final transformer layer is a strong default. \revised{In our
    benchmarks, multi-scale fusion provided no practical benefit over single-layer
    embeddings for either model under leakage-free evaluation}, and 1280 dimensions
    sufficed for both linear and nonlinear readouts.
    This reduces
    storage and computational costs by a factor of 5--10 compared to the common
    practice of concatenating 6--8 layers (up to 10,240 dimensions).

    \item \textbf{Standard PLMs (ESM-2, ProtBERT, etc.):} \revised{Under leakage-free
    evaluation, concatenation of 3--5 strategically spaced layers offers at most a
    marginal linear-readout gain ($+0.018$ in our experiments, from $\rho = 0.463$
    to 0.481). Switching the readout matters more than fusing layers: last-layer
    LightGBM reached $\rho = 0.508$, the best pooling configuration overall,
    achieving the benefits claimed for multi-scale fusion without the dimensionality
    cost.}
\end{enumerate}

Table~\ref{tab:cost} quantifies the practical trade-offs. Concat~5L uses
$4.9\times$ more features than single-layer approaches and incurs roughly
$25\times$ the computational cost for Ridge regression (the covariance matrix
grows from $1{,}322^2$ to $6{,}442^2$ entries), yet offered no benefit on
\revised{either model} in our benchmark \revised{and was matched or exceeded on
ESM-2 by switching from Ridge to LightGBM on a single layer}. Storing Concat~5L embeddings for a typical DMS
benchmark ($\sim$10,000 variants) requires 258~MB versus 53~MB for a single
layer, a gap that widens linearly with dataset size and becomes substantial
for genome-scale studies.

\begin{table}[ht]
\centering
\caption{Computational cost comparison across pooling strategies.
``Per-mutation'' storage assumes float32 precision. Time estimates are
per-dataset averages on an NVIDIA RTX 5090, including embedding extraction
and readout training across 20 ProteinGym DMS subsets.
\revised{Best achievable $\rho$ from Table~\ref{tab:pooling} (mutant arm,
modulo split).}}
\label{tab:cost}
\begin{tabular}{lcccc}
\toprule
& Last Layer & Mean L20--33 & Attention & Concat 5L \\
\midrule
Feature dimension                  & 1,322 & 1,322 & 1,322 & 6,442 \\
Storage per mutation (KB)          & 5.2   & 5.2   & 5.2   & 25.2 \\
Ridge design matrix (5,000 vars)   & $6.6 \times 10^6$ & same & same & $3.2 \times 10^7$ \\
Ridge covariance matrix entries    & $1.7 \times 10^6$ & same & same & $4.1 \times 10^7$ \\
Avg.\ time per dataset (min)       & $\sim$3 & $\sim$3 & $\sim$3 & $\sim$3 \\
\midrule
\multicolumn{5}{l}{\small \textbf{Best achievable $\rho$ (ESM-2 / SaProt):}} \\
\quad With Ridge                    & \revised{0.463 / 0.432} & \revised{0.473 / 0.433} & \revised{0.439 / 0.429} & \revised{0.481 / 0.434} \\
\quad With LightGBM                 & \textbf{\revised{0.508 / 0.448}} & \revised{0.457 / 0.439} & \revised{0.463 / 0.433} & \revised{0.495 / 0.436} \\
\bottomrule
\end{tabular}
\end{table}

In both cases, computational resources may be better directed toward feature types
that are plausibly complementary to PLM embeddings:
\begin{enumerate}
    \item Structural features (e.g., AlphaFold-predicted pLDDT, PAE~\cite{jumper2021alphafold},
          or Foldseek 3Di states~\cite{vankempen2024foldseek}), which complement
          sequence-based embeddings;
    \item Sequence-level representations that capture global protein properties
          beyond per-residue features;
    \item Multiple sequence alignment (MSA)-based approaches~\cite{riesselman2018deepsequence,
          hopf2017mutation} that leverage evolutionary information beyond what PLMs encode.
\end{enumerate}

As discussed under Boundary Conditions (Section~3.5), \revised{the position-lookup
control shows that} standard evaluation protocols can conflate mutation-effect learning with
position-level memorization. We recommend that future DMS benchmarks adopt
position-aware evaluation as a standard complement to random splitting.

\subsection{Limitations}

Several limitations should be noted. First, we tested the 650M-parameter
ESM-2 variant and its structure-aware derivative SaProt. Smaller (35M, 150M)
and larger (3B, 15B) variants may show different degrees of layer specialization,
and this remains to be tested. Second, our
experiments are restricted to per-residue mutation effect prediction; for other
tasks such as kinase function prediction~\cite{kumar2025}, subcellular
localization, or protein--protein interaction scoring~\cite{elshaffei2025},
layer-wise differences and the benefits of multi-scale fusion may be more
pronounced. Third, our pooling benchmark used a 20-dataset validation subset,
which is smaller than the \revised{63}-dataset probing analysis. \revised{Fourth,
the SaProt pooling runs use structure-masked inputs, whereas the probing and logit
lens runs include real 3Di inputs, and one dataset (influenza nucleoprotein) lacks
an AFDB structure. Fifth, our}
feature ablation (embedding-only vs.\ auxiliary-only) indicates that learned embeddings
dominate predictive performance in this setting, but it does
not rule out the possibility that nonlinear combinations of auxiliary features
could provide stronger baselines.
\revised{Sixth,} the
maximum sample size of 10,000 mutations per dataset (subsampled to 5,000 for
regression\revised{; 11 of 63 datasets were capped}) may limit sensitivity to
subtle layer-wise effects on the largest DMS screens ($>$200,000 mutations),
though the consistency of invariance scores across dataset-size categories argues
against this concern. \revised{Seventh,} our linear probing approach can only detect
information that is linearly separable; nonlinear interactions between layers
could theoretically exist but would require a different experimental design to
identify. \revised{Eighth}, position-level data splitting (Section~3.5) shows that
random mutation splits permit position-level leakage, \revised{which the lookup
control quantifies and which now motivates our reporting of leakage-free splits
throughout; future DMS benchmarks should adopt position-aware evaluation as a
standard protocol alongside random splitting. Finally, our equivalence margin}
($\Delta\rho = 0.01$) is \revised{a reporting convention rather than a
task-derived requirement, and} our
statistical analysis supports practical invariance by small effect sizes and
relative variation, but it does not constitute a pre-specified formal equivalence
test.

\subsection{Conclusion}

\revised{Under leakage-free evaluation, protein language model layers are close to
interchangeable for mutation effect prediction, with a small consistent final-layer
advantage; under the standard random split, apparent interchangeability is dominated
by position memorization.} Under standard random-split probing,
SaProt and ESM-2 showed near-flat per-layer performance despite large changes in
representational geometry. \revised{Multi-scale pooling added no practical benefit
over the final layer for either model under the unified leakage-free pipeline.
Structure tokens helped SaProt, but only under leakage-free evaluation.
Position-level evaluation showed that random splits inflate both absolute
performance and apparent invariance.} These findings argue for a more empirical use
of PLM layers: choose the layer strategy jointly with the readout model and the
evaluation protocol, and treat claims of layer-specific representations in PLMs with
the same scrutiny that has become standard in the NLP community~\cite{rogers2021bertology}.

%=============================================================================
\section*{Acknowledgments}
We thank the International
College of Pharmaceutical Innovation, SUDA-RCSI for institutional support.

\section*{Funding}
\revised{This work was supported by the Jiangsu Tepin Foundation (Z241922) and
the Health Research Board (ILP-PHR-2024-014).}

\section*{Competing Interests}
\revised{[AUTHOR\_INPUT\_NEEDED: competing interests declaration for all authors.]}

\section*{Data Availability}
Code and results are available at \url{https://github.com/schwarzeteeguo-creator/esm2-layer-invariance}.
\revised{An archived snapshot of the code, per-dataset results, and extracted
structure tokens for this revision is deposited at Zenodo, DOI:
[AUTHOR\_INPUT\_NEEDED].}
Raw ProteinGym data: \url{https://github.com/OATML-Markslab/ProteinGym}.

\newpage
\begin{thebibliography}{99}

% === References numbered by order of first citation in the text ===

\bibitem{rives2021biological}
Rives, A.\ \textit{et al.} (2021)
Biological structure and function emerge from scaling unsupervised learning to 250
million protein sequences.
\textit{Proceedings of the National Academy of Sciences}, \textbf{118}(15), e2016239118.

\bibitem{meier2021language}
Meier, J.\ \textit{et al.} (2021)
Language models enable zero-shot prediction of the effects of mutations on protein
function.
\textit{Advances in Neural Information Processing Systems}, \textbf{34}, 29287--29303.

\bibitem{lin2023evolutionary}
Lin, Z.\ \textit{et al.} (2023)
Evolutionary-scale prediction of atomic-level protein structure with a language model.
\textit{Science}, \textbf{379}(6637), 1123--1130.

\bibitem{frazer2021disease}
Frazer, J.\ \textit{et al.} (2021)
Disease variant prediction with deep generative models of evolutionary data.
\textit{Nature}, \textbf{599}(7883), 91--95.

\bibitem{hie2021learning}
Hie, B.\ \textit{et al.} (2021)
Learning the language of viral evolution and escape.
\textit{Science}, \textbf{371}(6526), 284--288.

\bibitem{elnaggar2022prottrans}
Elnaggar, A.\ \textit{et al.} (2022)
ProtTrans: Toward understanding the language of life through self-supervised learning.
\textit{IEEE Transactions on Pattern Analysis and Machine Intelligence}, \textbf{44}(10),
7112--7127.

\bibitem{brandes2022proteinbert}
Brandes, N.\ \textit{et al.} (2022)
ProteinBERT: A universal deep-learning model of protein sequence and function.
\textit{Bioinformatics}, \textbf{38}(8), 2102--2110.

\bibitem{rogers2021bertology}
Rogers, A., Kovaleva, O.\ and Rumshisky, A. (\revised{2020})
A primer in BERTology: What we know about how BERT works.
\textit{Transactions of the Association for Computational Linguistics}, \textbf{8},
842--866.

\bibitem{tenney2019bert}
Tenney, I., Das, D.\ and Pavlick, E. (2019)
BERT rediscovers the classical NLP pipeline.
\textit{Proceedings of the 57th Annual Meeting of the Association for Computational
Linguistics}, 4593--4601.

\bibitem{nostalgebraist2020}
nostalgebraist (2020)
Interpreting GPT: the logit lens.
\url{https://www.lesswrong.com/posts/AcKRB8wDpdaN6v6ru/}.

\bibitem{kumar2025}
Kumar, A.\ and Jha, I.P. (2025)
Layer probing improves kinase functional prediction with protein language models.
\textit{arXiv}, 2512.00376.

\bibitem{detlefsen2022}
Detlefsen, N.S.\ \textit{et al.} (2022)
Learning meaningful representations of protein sequences.
\textit{Nature Communications}, \textbf{13}(1), 1914.

\bibitem{elshaffei2025}
Alsamkary, H., Elshaffei, M.\ \textit{et al.} (2025)
Beyond simple concatenation: Fairly assessing PLM architectures for multi-chain
protein--protein interactions prediction.
\textit{arXiv}, 2505.20036.

\bibitem{rao2019tape}
Rao, R.\ \textit{et al.} (2019)
Evaluating protein transfer learning with TAPE.
\textit{Advances in Neural Information Processing Systems}, \textbf{32}.

\bibitem{notin2024proteingym}
Notin, P.\ \textit{et al.} (2024)
ProteinGym: Large-scale benchmarks for protein fitness prediction and design.
\textit{Advances in Neural Information Processing Systems}, \textbf{36}.

\bibitem{su2024saprot}
Su, J.\ \textit{et al.} (2024)
SaProt: Protein language modeling with structure-aware vocabulary.
\textit{International Conference on Learning Representations (ICLR)}.

\bibitem{vankempen2024foldseek}
van Kempen, M.\ \textit{et al.} (2024)
Fast and accurate protein structure search with Foldseek.
\textit{Nature Biotechnology}, \textbf{42}(2), 243--246.

\bibitem{roberts2017crossvalidation}
Roberts, D.R.\ \textit{et al.} (2017)
Cross-validation strategies for data with temporal, spatial, hierarchical, or
phylogenetic structure.
\textit{Ecography}, \textbf{40}(8), 913--929.

\bibitem{he2016deep}
He, K.\ \textit{et al.} (2016)
Deep residual learning for image recognition.
\textit{Proceedings of the IEEE Conference on Computer Vision and Pattern Recognition (CVPR)},
770--778.

\bibitem{jumper2021alphafold}
Jumper, J.\ \textit{et al.} (2021)
Highly accurate protein structure prediction with AlphaFold.
\textit{Nature}, \textbf{596}(7873), 583--589.

\bibitem{riesselman2018deepsequence}
Riesselman, A.J., Ingraham, J.B.\ and Marks, D.S. (2018)
Deep generative models of genetic variation capture the effects of mutations.
\textit{Nature Methods}, \textbf{15}(10), 816--822.

\bibitem{hopf2017mutation}
Hopf, T.A.\ \textit{et al.} (2017)
Mutation effects predicted from sequence co-variation.
\textit{Nature Biotechnology}, \textbf{35}(2), 128--135.

\bibitem{golub1979generalized}
\revised{Golub, G.H., Heath, M.\ and Wahba, G. (1979)
Generalized cross-validation as a method for choosing a good ridge parameter.
\textit{Technometrics}, \textbf{21}(2), 215--223.}

\bibitem{kornblith2019similarity}
\revised{Kornblith, S., Norouzi, M., Lee, H.\ and Hinton, G. (2019)
Similarity of neural network representations revisited.
\textit{Proceedings of the 36th International Conference on Machine Learning},
3519--3529.}

\bibitem{belrose2023eliciting}
\revised{Belrose, N. \textit{et al.} (2023)
Eliciting latent predictions from transformers with the tuned lens.
\textit{arXiv}, 2303.08112.}

\end{thebibliography}

\end{document}
